% -*- coding: utf-8 -*-
% 高数基础知识.tex
% 高数基础知识
\documentclass[UTF8]{ctexart}
\usepackage{geometry}

\geometry{a4paper,scale=0.8}

\usepackage{amsmath}
\usepackage{xcolor}
\usepackage{amssymb}
\usepackage{amsfonts}
\usepackage{enumitem}

% ===== 树状图/分类结构 =====
\usepackage{tikz}
\usetikzlibrary{trees,positioning}


\title{高数基础知识}
\author{埃及猪肉}
\date{\today}

\begin{document}
\maketitle
\tableofcontents


\section*{基础概念}

% 使用 cases 环境实现横向大括号树状图
\[
\begin{aligned}
&\text{解析式} \begin{cases}
    \text{代数式} \begin{cases}
        \text{有理式} \begin{cases}
            \text{整式} \begin{cases}
                \text{单项式} \\
                \text{多项式}
            \end{cases} \\
            \text{分式} \begin{cases}
                \text{真分式} \\
                \text{假分式}
            \end{cases}
        \end{cases} \\
        \text{无理式}
    \end{cases} \\
    \text{超越式}
\end{cases}
\end{aligned}
\]

\vspace{1em}

% 注释说明
\begin{enumerate}[label=\textcircled{\arabic*}]
    \item \textbf{有理式和无理式的区分}
    
    含有\textbf{字母}的根式运算的代数式是无理式。
    
    $\sqrt{2x}$ 是无理式，$\sqrt{2}x$ 是有理式。
    
    \item 含有字母的指数为无理数的指数运算、对数运算、三角运算和反三角运算的解析式，如
    \[
    \ln(x^2+y^2),\quad \arctan x,\quad a^{\sqrt{x}} \text{ 等}
    \]
\end{enumerate}

\end{document}